\section{A finite polar classification of detector spaces}
\label{sec:v85-polar}
The functional criterion admits a finite algebraic classification in a
natural category larger than polynomial and logarithmic detectors. Let
$J_0$ be an open clock interval, let $I$ be an open interval strictly to
its left, and let $\cF\subset C^1(J_0)$ be finite dimensional and contain
constants. In this section every derivative of an element of $\cF$ is
a rational function without poles on $J_0$. Put
\[
 E=\{e':e\in\cF\}\subset\R(T).
\]
Let $S$ be the finite set of poles occurring in $E$ and $S_I=S\cap I$.
Nonreal poles are understood in conjugate pairs. Partial fractions
identify $E$ with a subspace of a finite real coefficient space,
including its polynomial part. Thus membership and intersection tests
in $E$ are finite linear algebra on a specified basis.

For $c\in\Z^{S_I}$ write $c_+=(\max(c_s,0))_s$ and set
\begin{equation}\label{eq:v85-charge-set}
 \mathcal Z_2(S_I)=\left\{c\ne0:\ \sum_s c_s=0,\quad
                   \sum_s(c_s)_+\le2\right\},\qquad
 R_c(T)=\sum_{s\in S_I}\frac{c_s}{T-s}.
\end{equation}
This is a finite set of integer vectors: at most four entries are
nonzero, and each entry has absolute value at most two. For
$A\subset S_I$ put $D_A=\operatorname{span}\{(T-s)^{-2}:s\in A\}$.

\begin{theorem}[Polar classification]\label{thm:v85-polar}
The detector space $\cF$ is regularly separating on $I$ if and only if
\begin{align}
 R_c&\notin E &&\text{for every }c\in\mathcal Z_2(S_I),
                                      \label{eq:v85-polar-secant}\\
 E\cap D_A&=\{0\} &&\text{for every }A\subset S_I,
                          \quad 1\le |A|\le2.\label{eq:v85-polar-tangent}
\end{align}
Consequently regular admissibility is decidable by finitely many
coefficient-space membership and rank tests, with no search over
unknown action pairs. The first test classifies exact identification;
the second classifies differentiable local inversion.
\end{theorem}
\begin{proof}
Differentiating a nontrivial secant difference gives the rational
function with signed residue measure
\[
             \delta_x-\delta_y-\delta_{x'}+\delta_{y'}.
\]
If it belongs to $E$, every noncancelled pole belongs to $S_I$.
Its residue vector, when nonzero, belongs to
$\mathcal Z_2(S_I)$. If the vector is zero, equality of the signed
measures $\delta_x-\delta_y=\delta_{x'}-\delta_{y'}$, with both
ordered pairs off the diagonal, implies $x=x'$ and $y=y'$.
Since constants belong to $\cF$, membership of the derivative in $E$
is equivalent to membership of the original secant difference in
$\cF$: subtract an antiderivative in $\cF$, leaving a constant.

Conversely, every vector in \eqref{eq:v85-charge-set} is a secant
residue vector. If its positive mass is two, list its positive atoms,
with multiplicity, as $p_1,p_2$ and its negative atoms as $n_1,n_2$.
Take $(x,y)=(p_1,n_1)$ and $(x',y')=(n_2,p_2)$. Both pairs are
nondegenerate because positive and negative supports are disjoint.
If its positive mass is one, write $c=\delta_p-\delta_n$ and choose
$z\in I\setminus\{p,n\}$. The pairs $(p,z)$ and $(n,z)$ realize
$c$. This proves \eqref{eq:v85-polar-secant} in both directions.

The derivative in $T$ of a quotient tangent is
\[
       \left(-\frac{u}{T-x}+\frac{v}{T-y}\right)'
                 =\frac{u}{(T-x)^2}-\frac{v}{(T-y)^2}.
\]
Distinct double poles cannot cancel. Every pole with nonzero
coefficient must therefore lie in $S_I$ if this function belongs to
$E$. Conversely any nonzero element of $E\cap D_A$, $|A|\le2$, is
such a derivative; in the one-pole case choose an arbitrary different
second action with coefficient zero. Again constants remove the
integration ambiguity. These observations prove
\eqref{eq:v85-polar-tangent} and the theorem.
\end{proof}

The tests concern subspace membership, not only the presence of a
pole somewhere in a basis. For instance a coupled simple-pole direction
can fail the integer-residue test even when each individual coordinate
is absent. This distinction is essential for nonsaturated spaces.
No sign condition on a sampled annihilator occurs in the classification.

\subsection{Maximal spaces under principal-part closure}
Fix an ambient finite-dimensional rational space with prescribed pole
orders and prescribed polynomial degree. Call a subspace
\emph{principal-part saturated} if it is invariant under projection to
each individual real partial-fraction coordinate, to each conjugate
real block for nonreal poles, and to the polynomial part. In the
polynomial part no admissibility restriction is necessary. The next
statement is within this fixed ambient space; it does not posit a
largest finite-dimensional detector space without an ambient bound.

\begin{corollary}[Maximal saturated enlargements]
\label{cor:v85-maximal}
For a principal-part saturated derivative space $E$, regular separation
is equivalent to the following two conditions: no double-pole direction
$(T-s)^{-2}$ with $s\in I$ is present, and simple-pole directions
$(T-s)^{-1}$ are present at at most one point of $I$.
Every polynomial direction, every principal part at a pole outside
$I$, and every direction of order at least three at a pole in $I$
is admissible.

Among saturated subspaces of the fixed ambient space, the maximal
regularly separating ones contain all these admissible directions,
contain no forbidden double-pole direction, and contain exactly one
of the available simple-pole directions in $I$ when any is available.
If none is available there is one maximal saturated subspace. Detector
spaces are obtained by adjoining constants to antiderivatives.
\end{corollary}
\begin{proof}
Saturation makes membership coordinatewise. Two available simple-pole
directions give $R_c$ for $c=\delta_s-\delta_t$, whereas a single
simple-pole coordinate cannot support a nonzero vector of total
residue zero. A double-pole coordinate violates the tangent test
already with $|A|=1$. All other coordinates are disjoint from the
supports of the two obstruction families. Theorem~\ref{thm:v85-polar}
therefore gives the first assertion. Adding any missing safe coordinate
preserves both tests. Omitting all available simple-pole coordinates
is not maximal; including two is inadmissible. This proves the
maximality statement, including all alternatives.
\end{proof}
Thus the closure operation has a precise answer. In particular a
third-order pole in the derivative, whose primitive is a second-order
rational detector, need not destroy regularity; the forbidden order
for regular inversion is specifically two. Regularity and mere
set-theoretic identification have different polar boundaries.

\subsection{One pole budget for uniform clock designs}
Let $Q_E$ be the monic least common denominator of $E$, with degree
$D$, and take $Q_E=1$ if $E$ consists of polynomials. Let $\ell(E)$
be the largest order at infinity of a nonzero member of $E$: a
function with leading term $aT^j$ has order $j$. Put
\[
                  b(E)=\max\{\ell(E),-2\},
\]
with $\ell(\{0\})=-\infty$. Both $D$ and $b(E)$ are read from the
same coefficient space as the polar tests.

\begin{theorem}[Polar clock budget]\label{thm:v85-budget}
If the tests of Theorem~\ref{thm:v85-polar} hold, then any
\begin{equation}\label{eq:v85-budget}
                         N=D+6+b(E)
\end{equation}
distinct clocks in $J_0$ identify $e\in\cF$ and the ordered unequal
actions, and the sampled map is an immersion. Its inverse is locally
smooth and uniformly Lipschitz on each compact separated action class,
in the sense of \eqref{eq:v84-design}. The count is sufficient, not
in general minimal.
\end{theorem}
\begin{proof}
For either a scalar collision or a scalar tangent, call the associated
function $H$. Its derivative is the sum of a member of $E$ and a
rational function with at most four moving pole factors, counted with
multiplicity, which is $O(T^{-2})$ at infinity. Multiply $H'$ by
$Q_E$ and those four factors, omitting unnecessary factors if fewer
are present. The result is a polynomial of degree at most
$D+4+b(E)$. If $H$ vanishes at the $N$ clocks, Rolle's theorem gives
$N-1=D+5+b(E)$ distinct zeros of $H'$ in the clock interval, where
the denominator never vanishes. The numerator is zero. Thus $H$ is
constant, and a sampled zero makes it identically zero. The secant
test then gives identical action pairs and zero nuisance difference;
the tangent test gives a zero parameter tangent. This proves global
injectivity and full differential rank. Smooth local inversion follows
from a nonsingular coordinate minor. The normalized-difference and
normalized-tangent compactness proof of Theorem~\ref{thm:v84-design}
applies to this fixed clock set and proves its compact-class estimate.
\end{proof}
For protected logarithms and polynomial log-detectors of degree $q$,
$D=k$ and $b(E)=q-1$ for $q\ge1$, while $b(E)=-1$ for $q=0$ and
$k>0$. This recovers the $q+k+5$ sufficient count. For
$\cF=\operatorname{span}\{1,\log((T-a)/(T-b))\}$, $D=2$ and
$b(E)=-2$, giving six clocks. The improved count is the cancellation
of the total residue at infinity, not positivity of an annihilator.
The sharp polynomial count in Theorem~\ref{thm:v83-scalar} is stronger
than \eqref{eq:v85-budget}; it uses the additional oriented-interval
structure. For several action thresholds, the identical argument uses
$4n$ moving factors and the integer-residue balance underlying
Theorem~\ref{thm:v84-divisors}. A common action divisor is then a
projective gauge, rather than a failure of the partial-fraction test.

\subsection{Quantitative conventions and comparison}
Throughout the function-space design statements, $|a-a'|$ denotes
the Euclidean norm on $\R^2$. A compact $K\subset\mathcal A$ has
$\inf_K|x-y|>0$ and positive distance from the clock singularities.
For rational-derivative spaces, compact conditioning constants also
depend on the fixed coefficient basis and on the separation of moving
actions from the fixed poles when such separations enter that class.
An exact residue identity is not a uniform conditioning assertion as
these quantities degenerate.

The closest general predecessors to the finite-design and analytic
incidence arguments are the analytic experiment-count theory of
Sontag~\cite{V85Sontag} and finite-measurement stability results of
Alberti--Santacesaria~\cite{V85AS}. The collision incidence here has
$d$ nuisance and four action coordinates; its dimension is $d+4$.
The normalized tangent family has $d+1$ coefficient coordinates and
two action coordinates, hence dimension $d+3$. These dimensions explain
the sufficient generic count, not a random-design condition number
or a new general embedding principle. The new structural conclusion
in this section is instead the finite charge and double-pole
classification, including maximal spaces under principal-part closure.
Partial fractions and rational zero counting are classical tools;
Prony-type recovery~\cite{V84BY} uses related singular-support data.
Here the obstruction vectors classify which unknown detector
subspaces preserve the nonlinear action quotient before a clock design
is chosen.
